\section{A constructed ordinary Eisenstein seed}
\label{sec:constructed-eisenstein-seed}

We construct the classical seed used by the rank-three branch.
The construction below produces actual cuspidal eigenforms.
Its last paragraph distinguishes this existence statement from
placing the seed in the required localized unitary support.

Let \(\ell\ge5\), fix an embedding of the algebraic coefficient
fields into \(\overline{\mathbb Q}_\ell\), and write
\(\mathfrak m_\ell\) for the ideal of elements of positive
\(\ell\)-adic valuation.
Take a primitive nontrivial Dirichlet character \(\varphi\)
of conductor \(M>1\), with \(\ell\nmid M\).
Choose a prime \(w\nmid M\) and a nontrivial character
\(\theta\) of conductor \(w\) such that
\[
       w\equiv1\pmod\ell,\qquad
       \varphi(w)=1,\qquad \theta^\ell=1 .
 \tag{ES1}
\]
Set \(R=Mw\) and \(\chi=\varphi\theta\), interpreted at level
\(R\) by the Chinese remainder theorem.
It is primitive: the two conductors are coprime and each
factor has its stated conductor.
In particular \(\ell\nmid R\).

We shall prove that, for every sufficiently large \(k\) with
\(\chi(-1)=(-1)^k\), there is a normalized cuspidal eigenform
\(f\) of weight \(k\), level \(R\), and character \(\chi\), with
\[
          a_p(f)\equiv1+\chi(p)p^{k-1}\pmod{\mathfrak m_\ell}
                       \quad\hbox{for every prime }p .
 \tag{ES2}
\]
This includes \(p\mid R\) and \(p=\ell\).
The latter congruence makes \(f\) ordinary.

\subsection{Integral forms with a prescribed constant term}
For \(r\ge1\), let \(P_r\in\mathbb Z[X]\) be defined by
\[
 P_1=X,\qquad
 P_{r+1}=X((1-X)P'_r+rP_r),\qquad
                    L_r(x)=\frac{P_r(x)}{(1-x)^r}.
\]
For \(|x|<1\), this is
\(\sum_{n\ge1}n^{r-1}x^n\); differentiation proves the recursion.
In particular
\[
 L_3(x)=\frac{x(1+x)}{(1-x)^3},\qquad
 L_5(x)=L_3(x)\left(1+12\frac{x}{(1-x)^2}\right).
 \tag{ES3}
\]
If \(x\ne1\) has finite order prime to \(\ell\), then
\(1-x\) is an \(\ell\)-adic unit.
For example, reduction preserves all roots of unity of
prime-to-\(\ell\) order: the geometric-sum identity shows
that two such roots with equal residue are equal.
Thus every \(L_r(x)\) is integral.
If also \(x^2\ne1\), then \(L_3(x)\) is a unit.

We record a separation calculation for these constant terms.
If \(x,y\) are \(R\)-th roots, both primitive, and
\(y\ne x,x^{-1}\), put
\[
                    u=\frac{x}{(1-x)^2},\qquad
                    v=\frac{y}{(1-y)^2}.
\]
Both are units, and
\[
 u-v=\frac{x(1-y/x)(1-xy)}{(1-x)^2(1-y)^2}
\]
is a unit as well.
Let \(A=L_3(x)^2\), \(B=L_3(y)^2\).
If \(A-B\) is a unit it already separates \(x,y\).
Otherwise (ES3) gives
\[
 L_3(x)L_5(x)-L_3(y)L_5(y)
       =12A(u-v)+(A-B)(1+12v).
\]
Its first term is a unit and its second term lies in
\(\mathfrak m_\ell\), since \(\ell\ge5\).
Their sum is a unit.
Consequently one of the constant-term functions
\(L_3^2\), \(L_3L_5\), of weights \(6,8\), separates any
two primitive roots not equal up to inversion.
This is the generic-prime specialization of the source's
weight-\(4,6,8\) calculation.

Here are the actual modular forms to which we apply it.
For \(a\in(\mathbb Z/R)^\times\) and \(r\ge3\), set
\[
 {\cal E}_{r,a}(z)=
 \frac{R^r}{(-2\pi i)^r/(r-1)!}
       \sum_{\substack{(c,d)\in\mathbb Z^2\\
                 c\equiv0\ (R),\ d\equiv a\ (R)}}
                                    (cz+d)^{-r}.
 \tag{ES4}
\]
The sum is absolutely and locally uniformly convergent.
Changing the summation variables proves its
\(\Gamma_1(R)\)-modularity, and for
\(\gamma\in\Gamma_0(R)\),
\[
             {\cal E}_{r,a}|_r\gamma
                           ={\cal E}_{r,a\gamma_{22}}.
 \tag{ES5}
\]
The same convergence estimate after a cusp change proves
holomorphy at all cusps.
If \(R\nmid\gamma_{21}\), the constant term of
\({\cal E}_{r,a}|_r\gamma\) is zero.
Indeed a constant term would require a lattice vector whose
transformed first coordinate is zero; the two congruences in
(ES4), with \(a\) a unit, would then force
\(R\mid\gamma_{21}\).
Equivalently these forms vanish at every cusp not equivalent
to infinity under \(\Gamma_0(R)\).

Writing \(\zeta=e^{2\pi i/R}\), the elementary Lipschitz
expansion of the inner sum gives
\[
 \begin{split}
 a_0({\cal E}_{r,a})&=L_r(\zeta^a),\\
 a_n({\cal E}_{r,a})&=
   \sum_{d\mid n}d^{r-1}
       \bigl(\zeta^{ad}+(-1)^r\zeta^{-ad}\bigr),\qquad n\ge1 .
 \end{split}
 \tag{ES6}
\]
One obtains this by summing first over the second lattice
coordinate, expanding the resulting geometric series for
positive first coordinate, and pairing negative first
coordinates with positive ones.
Absolute convergence for \(r>2\) justifies these operations.
These formulas prove integrality of every Fourier coefficient.
They also prove stability under every automorphism of
\(\mathbb C\): such an automorphism replaces \(\zeta\) by
\(\zeta^s\), and hence replaces the family by
\({\cal E}_{r,sa}\).
No invariance of the transcendental normalizing factor in
(ES4) is asserted; (ES6) is the coefficient calculation
which proves the required conjugation statement.

For every \(d\in(\mathbb Z/R)^\times\setminus\{1,-1\}\), use
the separation calculation to choose
\[
 H_a={\cal E}_{3,a}^2\quad\hbox{or}\quad
 H_a={\cal E}_{3,a}{\cal E}_{5,a}
\]
whose constant terms at \(a=1,d\) differ by a unit.
Let \(r_d=6\) or \(8\), accordingly.
There is an integral level-one form \(J_{r_d}\) with constant
term one: take \(E_6\) or \(E_4^2\).
The family
\[
                 H_a-a_0(H_d)J_{r_d}
 \tag{ES7}
\]
has constant term zero at \(a=d\) and a unit at \(a=1\).
It is integral and has the family transformation law (ES5).
Take the product of (ES7) over all \(d\ne\pm1\).
It is an even-weight family, of some weight \(m_0\), whose
constant term is zero except possibly at \(a=\pm1\),
and is a unit at \(a=1\).
All its coefficients lie in a fixed cyclotomic number field.
One may take \(m_0\le8(\phi(R)-2)\).

To reach every sufficiently large weight, use
\({\cal E}_{3,a}\), whose constant at \(a=1\) is a unit.
For odd \(j\ge10\), multiply it by an integral level-one
form of weight \(j-3\) and constant one.
For even \(j\ge10\), use \({\cal E}_{3,a}^2\) and a level-one
form of weight \(j-6\).
Such level-one forms exist in every nonnegative even weight
other than \(2\), as products of \(E_4,E_6\).
Multiplying by the previous product gives, for every
\(k\ge m_0+10\), an integral weight-\(k\) family \(F_a\)
which vanishes at the other cusps, satisfies (ES5), and has
\[
  u=a_0(F_1)\in\mathcal O_\ell^\times,\qquad
  a_0(F_a)=0\quad(a\ne\pm1),\qquad
                   a_0(F_{-1})=(-1)^ku .
\]
For any character \(\chi\) with the weight parity, take
\[
             G_\chi=\frac1{2u}
               \sum_{a\in(\mathbb Z/R)^\times}\chi(a^{-1})F_a .
 \tag{ES8}
\]
The transformation law makes its character \(\chi\).
It is integral, has constant term one, and vanishes at all
the other cusps.
Only the unit \(2u\) was divided out.
In particular this projection remains integral even if
\(\ell\mid\phi(R)\).
The coefficient formulas and the construction give modular
conjugates of \(G_\chi\) for every automorphism of
\(\mathbb C\).

\subsection{The small Eisenstein constant, proved by distribution}
The character projection of (ES4), divided by its Gauss sum,
gives a form \(E_\chi\) of weight \(k\), character \(\chi\),
vanishing at the other cusps, with
\[
 \begin{split}
 a_n(E_\chi)&=\sum_{d\mid n}\chi(d)d^{k-1}\quad(n\ge1),\\
 2\tau(\chi^{-1})a_0(E_\chi)
        &=\sum_{a\bmod R}\chi^{-1}(a)L_k(\zeta^a),
            \qquad \tau(\chi^{-1})=\sum_{a\bmod R}
                                      \chi^{-1}(a)\zeta^a .
 \end{split}
 \tag{ES9}
\]
The usual Gauss identity
\(\tau(\chi)\tau(\chi^{-1})\chi(-1)=R\) shows that both
Gauss sums are \(\ell\)-adic units, since \(\ell\nmid R\).
This identity follows by applying the finite Fourier
transform twice to the primitive character.
All the nonconstant coefficients in (ES9) are integral.

We prove \(c=a_0(E_\chi)\in\mathfrak m_\ell\) directly.
For \(x^w\ne1\), the geometric-series expansion and then
differentiation give
\[
             \sum_{j=0}^{w-1}L_k(x\zeta_w^j)
                                       =w^kL_k(x^w).
 \tag{ES10}
\]
Both sides are rational functions, so equality initially
proved for \(|x|<1\) proves it wherever their denominators
are nonzero.
Let
\[
          T_\varphi=\sum_{b\bmod M}
                         \varphi^{-1}(b)L_k(\zeta^{wb}).
\]
This is integral, since all nonzero contributing roots have
order dividing \(M\), prime to \(\ell\).
Split the sum on the right of (ES9) into multiples of \(w\)
and their complement, and use (ES10) on the full sum with
the \(\theta\)-factor omitted.
Because \(\varphi(w)=1\), the result is the exact identity
\[
 \begin{split}
 \sum_{a\bmod R}\chi^{-1}(a)L_k(\zeta^a)
  =\sum_{\substack{0\le n<R\\w\nmid n}}
        (\theta^{-1}(n)-1)\varphi^{-1}(n)L_k(\zeta^n)\\
                 +(w^k-1)T_\varphi .
 \end{split}
 \tag{ES11}
\]
The terms with zero character value contribute zero and may
be omitted before evaluating any rational function at \(1\).
Every \(\ell\)-th root of unity is congruent to \(1\) modulo
\(\mathfrak m_\ell\), and \(w^k\equiv1\pmod\ell\).
Each term of (ES11) therefore lies in
\(\mathfrak m_\ell\).
Since \(2\tau(\chi^{-1})\) is a unit, (ES9) proves
\[
                               c\in\mathfrak m_\ell .
 \tag{ES12}
\]

Consequently
\[
                             h=E_\chi-cG_\chi
\]
is a cusp form: it vanishes at the other cusps, and its
constant term at infinity is zero.
The character law transports this vanishing to the remaining
cusps equivalent to infinity under \(\Gamma_0(R)\).
Its coefficients are integral and
\[
 a_n(h)\equiv\sum_{d\mid n}\chi(d)d^{k-1}
                    \pmod{\mathfrak m_\ell}\quad(n\ge1),
                   \qquad a_1(h)\equiv1 .
 \tag{ES13}
\]
Thus its residual form is nonzero.

\subsection{Lifting the simultaneous eigenvalues}
The divisor-sum identity at every prime, including those
dividing \(R\), is
\[
 \sigma_\chi(pn)+\chi(p)p^{k-1}
       {\boldsymbol1}_{p\mid n}\sigma_\chi(n/p)
        =(1+\chi(p)p^{k-1})\sigma_\chi(n).
\]
It follows either by splitting the divisors according to
their \(p\)-valuation or by their geometric-series factor.
Applying it to (ES13) makes \(h\) a simultaneous residual
eigenvector for all the prime Hecke operators, with the
eigenvalues in (ES2).

Here is the integral eigenvalue lifting argument.
Choose one finite coefficient field containing every
coefficient in the explicit construction of \(h\), and let
\(\mathcal O\) be its \(\ell\)-adic integer ring.
Let \(V\) be the span over its fraction field of all Hecke
translates of \(h\).
It is finite dimensional.
Indeed all these translates have coefficients in that one
field, and linear independence of finitely many such
\(q\)-series is detected by finitely many coefficients.
Their rank over the coefficient field equals their rank
over \(\mathbb C\); the complex cusp-form space is finite
dimensional.
The lattice of elements of \(V\) with all coefficients
in \(\mathcal O\) embeds in \(\mathcal O^B\) using a finite
injective list of coefficient evaluations.
It is therefore finite and free.
It spans \(V\), since the integral Hecke translates of \(h\)
span \(V\).
The Hecke operators preserve it by their coefficient formula.

Their image \(\mathbb T\) in its endomorphism ring is a
finite, torsion-free, commutative \(\mathcal O\)-algebra.
The nonzero residual eigenvector (ES13) gives a maximal
ideal \(\mathfrak n\) and the specified residual eigenvalues.
Choose a minimal prime of \(\mathbb T\) below \(\mathfrak n\).
It does not contain the uniformizer, since the latter is a
non-zero-divisor in \(\mathbb T\).
The resulting generic character takes values in a finite
extension of the coefficient field and has a place reducing
to \(\mathfrak n\).
It occurs on \(V\): the faithful finite module has support
all of \(\operatorname{Spec}\mathbb T\), and the socle of
the corresponding local Artinian generic algebra gives a
common eigenvector.
This is the elementary Deligne--Serre lifting argument.

For a nonzero common eigenvector the first coefficient is
nonzero.
Otherwise the Hecke recurrence and induction on \(n\), using
a prime divisor of \(n\), make every positive coefficient
zero; cuspidality makes the zeroth coefficient zero as well.
Normalize that first coefficient to one.
The resulting eigenform satisfies (ES2).
Since \(\chi\) has conductor exactly \(R\), it cannot be
old from a proper divisor of \(R\): the central character
of such an old constituent would have smaller conductor.
Thus the form is a newform of the stated primitive level.
This completes the construction of the ordinary cuspidal
Eisenstein congruence.

\subsection{Its Galois representation and its rank-three adjoint}
We specify the classical inputs at this point, so that their
fields and hypotheses are visible.
For a normalized holomorphic newform of weight \(k\ge3\)
over \(\mathbb Q\), Deligne's construction supplies its
continuous absolutely irreducible representation \(\rho_f\),
with
\[
 \operatorname{tr}\rho_f(\operatorname{Frob}_p)=a_p(f),
 \qquad
 \det\rho_f=\varepsilon^{k-1}\chi .
 \tag{ES14}
\]
Here we use arithmetic Frobenius and its corresponding
standard dual convention; replacing every representation
by the dual gives the inverse-Frobenius convention used
elsewhere in the source.
Its Hodge--Tate weights in the convention
\(\mathrm{HT}(\varepsilon)=1\) are \(0,k-1\).
The prime \(\ell\) does not divide its level, and (ES2)
makes \(a_\ell(f)\) a unit.
The ordinary local theorem therefore gives an upper
triangular restriction at \(\ell\), whose unramified
quotient takes Frobenius to the unit root of
\[
                    X^2-a_\ell(f)X+\chi(\ell)\ell^{k-1}.
\]
At primes different from \(\ell\), local--global compatibility
identifies its Weil--Deligne representation with that of
the newform.
These are applications to a classical newform over
\(\mathbb Q\), of weight at least three, prime-to-\(\ell\)
level, and ordinary unit root; no Hilbert even-degree or
residual irreducibility hypothesis is inserted.

Equations (ES2) and (ES14), followed by density of Frobenius
and character independence in dimension two, give
\[
                  \bar\rho_f^{\,\mathrm{ss}}
                        =1\oplus\overline{\varepsilon^{k-1}
                                                      \varphi}.
 \tag{ES15}
\]
The \(\theta\)-factor disappears modulo
\(\mathfrak m_\ell\).
At every \(p\mid R\), primitive conductor forces the local
representation to be a principal series with one
unramified character and the ramified character prescribed
by \(\chi\).
Indeed the conductor of the central character already
equals the entire newform conductor; a special or
supercuspidal constituent would have larger representation
conductor than central-character conductor.
Thus its inertia type is exactly \(1\oplus\chi|_{I_p}\)
and its monodromy operator is zero.

Put \(m=k-1\), and suppose, as in the principal seed data,
\[
                  \overline{\varepsilon^m\varphi}^{\,2}\ne1
                         \quad\hbox{on }G_{\mathbb Q}.
 \tag{ES16}
\]
Then \(\rho_f\) remains absolutely irreducible on every
open subgroup.
Here is a proof avoiding any image-size assertion.
If its restriction to an open normal subgroup were
reducible, prime-dimensional Clifford theory gives either
a scalar restriction or two lines interchanged transitively.
A scalar restriction would make its projective image finite.
On an open local subgroup at \(\ell\) its two Hodge--Tate
weights would then coincide, contrary to \(0,m\), \(m>0\).
Two interchanged lines give a quadratic self-twist
\(\alpha\ne1\).
Reducing that twist in (ES15), in odd characteristic, gives
\[
        \{1,\bar\Phi\}=\{\bar\alpha,\bar\alpha\bar\Phi\},
                         \qquad \Phi=\varepsilon^m\varphi.
\]
Since the nontrivial quadratic character stays nontrivial
modulo \(\ell\), these multisets force
\(\bar\alpha=\bar\Phi\) and \(\bar\Phi^2=1\), contradicting
(ES16).

In characteristic zero this also makes
\(\operatorname{ad}^0\rho_f\) irreducible on every open
subgroup.
The connected Zariski closure of \(\rho_f\) is reductive:
its unipotent radical has nonzero fixed vectors, and
normality and irreducibility make it act trivially.
It acts irreducibly in dimension two and cannot be a torus.
Its nonzero semisimple derived Lie algebra is therefore
\(\mathfrak{sl}_2\); the connected group contains
\(\mathrm{SL}_2\).
The adjoint action on its three-dimensional simple Lie
algebra is irreducible.
Every open subgroup has the same connected Zariski closure,
which proves the assertion.

Now let \(L/\mathbb Q\) be the finite abelian CM extension
in the seed packet, with
\[
          \eta|_{G_L}=\theta|_{G_L}=1,\qquad
                         \varphi=\chi'\eta .
\]
For a prescribed residue weight \(j\), choose
\[
                         m\equiv j\pmod{\ell-1}
\]
as large as desired and \(k=m+1\) above the constructed
threshold.
The seed parity is
\(\varphi(-1)=(-1)^{j+1}\), so it gives the required
weight parity.
On \(G_L\), (ES15) and the adjoint formula give the exact
rank-three residual semisimplification
\[
 \overline{\operatorname{ad}^0\rho_f|_{G_L}}^{\,\mathrm{ss}}
       =\bar\Psi\oplus1\oplus\bar\Psi^{-1},
                          \qquad \Psi=\varepsilon^j\chi' .
 \tag{ES17}
\]
The ordinary flag has Hodge--Tate weights \(m,0,-m\).
Its ordered characters are the ratio of the two ordinary
characters of \(\rho_f\), the trivial character, and the
inverse ratio.
Their reductions are exactly the ordered tuple in (ES17).

The trace form on trace-zero two-by-two matrices is
nondegenerate, symmetric, and invariant under conjugation.
For a complex conjugation \(c\), multiply its Gram matrix
by \(\operatorname{ad}^0\rho_f(c)\).
The result is symmetric, since \(c^2=1\), and gives
\[
       r(g)^{\mathsf t}P\,r(cgc^{-1})=P,
                  \qquad r=\operatorname{ad}^0\rho_f|_{G_L}.
 \tag{ES18}
\]
Thus \(r\varepsilon^{-1}\) has the standard unitary
multiplier \(\varepsilon^{-2}\).
Away from \(\ell\) its finite inertia type is
\(\chi'\oplus1\oplus\chi'^{-1}\); the \(\eta,\theta\)
factors disappear on \(G_L\).
In particular it is unramified away from the ramification
set allowed in the seed packet.
All these assertions concern the actual adjoint
representation, rather than a determinant with the same
generic roots.

\subsection{The automorphic transfer and local root bounds}
There is a direct classical realization of this seed.
Iterate cyclic base change for \(\mathrm{GL}_2\) along a
prime-degree tower for the abelian extension \(L/\mathbb Q\).
The base change is cuspidal: a noncuspidal step would be
the dihedral/self-twist case, making the attached
two-dimensional representation reducible on a finite-index
subgroup, already excluded above.
At each step the unramified parameters match restriction,
so the Galois realization is \(\rho_f|_{G_L}\).
It has no nontrivial self-twist, by the same open-subgroup
irreducibility.
The adjoint lift of Gelbart--Jacquet, Theorem~9.3, therefore
gives a cuspidal \(\mathrm{GL}_3(\mathbb A_L)\)-representation
with the local adjoint parameters.
It realizes \(r\), and after the indicated algebraic scalar
normalization realizes \(r\varepsilon^{-1}\).
The weights are regular, since \(m>0\), and (ES18) gives
the required conjugate self-duality.

The same theorem supplies the unlinked local property of this
seed without a new Selmer hypothesis.
At a finite place, unitarize the \(\mathrm{GL}_2\)-representation.
For a complementary principal-series component its two
Weil eigenvalue absolute values have ratio \(q^{2t}\).
Gelbart--Jacquet, Theorem~9.3(4), gives \(|t|<1/4\).
For a tempered component take \(t=0\).
The primitive-level ramified components above are principal
series with ramified ratio; their unitarity makes both
inducing characters unitary, so they too have \(t=0\).
For the adjoint representation, the three absolute values
are consequently
\[
                              q^{2t},\quad1,\quad q^{-2t}.
\]
The absolute value of a ratio between any two is strictly
between \(q^{-1}\) and \(q\).
After local base change of residue degree \(f\), the exponents
and the logarithm of the residue norm both multiply by \(f\),
so the same strict inequality holds with the new norm.
Finite inertial character factors have absolute value one.
Thus no two adjoint Frobenius roots have ratio equal to the
local residue norm or its inverse.
The monodromy is zero at the finite ramified places under
consideration.
This proves the actual unlinked condition for this seed;
it is a property of the constructed local parameters,
rather than a condition imposed on an arbitrary deformation.

The classical sources just used are Deligne,
\emph{Formes modulaires et repr\'esentations \(\ell\)-adiques},
Bourbaki~355, for the two-dimensional realization;
the ordinary local theorem in Wiles,
\emph{On ordinary \(\lambda\)-adic representations associated
to modular forms}, \emph{Invent.\ Math.}~94 (1988), 529--573;
Carayol, \emph{Sur les repr\'esentations \(\ell\)-adiques
associ\'ees aux formes modulaires de Hilbert}, Theorem~A,
whose base field here is \(\mathbb Q\), of odd degree;
Langlands, \emph{Base Change for \(\mathrm{GL}(2)\)},
Chapter~2, for the cyclic transfers;
and Gelbart--Jacquet, \emph{A relation between automorphic
representations of \(\mathrm{GL}(2)\) and \(\mathrm{GL}(3)\)},
Theorem~9.3, for the adjoint lift.
Primary copies are respectively
\url{https://archive.numdam.org/item/SB_1968-1969__11__139_0/},
\url{https://doi.org/10.1007/BF01394275},
\url{https://archive.numdam.org/articles/10.24033/asens.1512/},
\url{https://publications.ias.edu/sites/default/files/base-change-for-gl2_rpl.pdf},
and
\url{https://www.numdam.org/articles/10.24033/asens.1355/}.
The specific ordinary matrix statement, including the
unit-root quotient, is matched with the duality convention
in Lemma~\ref{lem:sr-stable-three-seed-insertion}.

The constructive part is the actual
\srcid{UnlinkedHeckeRoots.PropEC.propEC} producer.
Its ingredients at the frozen tag are
\srcid{IdeleGenerators.HolomorphicIrreducibility.EisSlash},
\srcid{EisVan}, \srcid{EisCoef},
\srcid{EulerianNumeratorThree.lkSep},
\srcid{LevelOne.levelOne}, and \srcid{ConstSmall.constSmall};
the conjugation-preserving version is
\srcid{UnlinkedHeckeRoots.UnitFamC.unitFamC},
\srcid{UnitFormC.unitFormC}, and \srcid{EigLiftC.eigLiftC}.
They are under
\srcid{Development.AutomorphicArithmetic.AutomorphicVarieties}.
The subsequent source calls are
\srcid{ThresholdArguments.DenseThresholdCases.WEGCore},
\srcid{TPropEGW}, \srcid{TSeedInW}, and
\srcid{AdjointLiftFromEisensteinSeedOfHyps}.

What has been produced is an actual ordinary classical
eigenform and an actual strongly irreducible polarized
rank-three automorphic adjoint with the required residual
and finite inertia packet, at arbitrarily large compatible
weights.
The field construction below chooses the definite unitary
pure inner form and retains the scalar Frobenius on the
selected tower. Lemma~\ref{lem:sr-stable-three-seed-insertion}
then supplies its ordinary local vector and localization
in the same Hecke module and determinant ring. That
insertion uses the actual characteristic-zero seed and
its local types, as well as its residual determinant.


\subsection{The Eisenstein seed at coefficient prime \(3\)}
\label{subsec:sr-eisenstein-three}
The constructed Eisenstein seed also works at \(\ell=3\)
on the actual prime-to-\(3\) level.
Only the constant-term separation calculation changes;
the absolutely convergent Eisenstein series have weight
at least three throughout.

Use the notation of
Section~\ref{sec:constructed-eisenstein-seed}.
For \(x\ne1\), put \(u=x/(1-x)^2\).
The Eulerian recurrence gives
\[
        L_4(x)=\frac{x(1+4x+x^2)}{(1-x)^4}
                                     =u+6u^2 .
 \tag{E3.1}
\]
If \(x,y\) are primitive roots of order \(R\), \(3\nmid R\),
and \(y\ne x,x^{-1}\), the same factorization as before gives
\[
 u-v=\frac{x(1-y/x)(1-xy)}
                         {(1-x)^2(1-y)^2}
                      \in\mathcal O_3^\times .
\]
All denominators and numerator factors are units, since
reduction preserves roots of unity of order prime to \(3\).
Consequently
\[
          L_4(x)-L_4(y)
               =(u-v)\bigl(1+6(u+v)\bigr)
                          \in\mathcal O_3^\times .
 \tag{E3.2}
\]
The second factor is congruent to \(1\) modulo \(3\).

Thus, for each \(d\ne\pm1\), replace the weight-\(6\)
or weight-\(8\) separating family in (ES7) by
\[
                H_a={\cal E}_{4,a},\qquad
                H_a-L_4(\zeta^d)E_4 .
 \tag{E3.3}
\]
The forms \({\cal E}_{4,a}\) are exactly the absolutely
convergent series (ES4); their constant terms and
coefficients are given by (ES6).
The integral level-one form \(E_4\) has constant term one.
Equation (E3.2) gives the required zero at \(a=d\)
and unit at \(a=1\).
The product has weight
\(m_0=4(\phi(R)-2)\).

Here \(R=Mw\), with \(M>1\) and a fresh prime
\(w\equiv1\pmod3\), so primitive \(R\)-th roots are not
\(1\) or \(-1\).
Their \(L_3\)-values are units.
Multiply this product by \({\cal E}_{3,a}\) in odd
remaining weight and by \({\cal E}_{3,a}^2\) in even
remaining weight, followed by a suitable product of
\(E_4,E_6\).
As in the original argument this reaches every
\(k\ge m_0+10\), with the chosen parity.
These final \({\cal E}_{3,a}\)-factors give vanishing
at every cusp not equivalent to infinity under
\(\Gamma_0(R)\).
The projection divides only by the unit \(2u\).
It therefore gives the integral \(G_\chi\) with
constant term one and the prescribed character.

The distribution identity (ES10)--(ES11) applies unchanged
to an auxiliary character \(\theta\) of exact order \(3\).
Its values are congruent to \(1\) modulo the maximal ideal;
also \(w^k-1\) is divisible by \(3\).
The Gauss sum is a unit because \(3\nmid R\).
Thus the Eisenstein constant is in the maximal ideal,
and \(E_\chi-cG_\chi\) is the same nonzero residual
cuspidal eigenform.
The finite integral Hecke-algebra lifting argument,
the primitive-level newform argument, Deligne realization
and ordinary local theorem impose no restriction
\(\ell\ge5\).
All of them now apply at \(3\).
In particular the all-odd-prime Eisenstein seed, its
rank-two normalization and its rank-three adjoint seed
are constructed at the prime required by the Gaussian
branch.
There is no weight-two regularization or inverted
factor \(3\) in this extension.


\section{The seed packet at every generic coefficient prime}
\label{sec:seed-all-generic-primes}

The prime-to-\(\ell\) hypothesis in the classical
Eisenstein construction concerns its finite character
conductor, not the original host conductor.
We make the necessary adjustment explicitly.
This also constructs a common CM field in which the
distinguished scalar Frobenius retains its required
cyclotomic value.

Let \(N=2^{h+2}qd\) be the principal host conductor, with
\(q,d>7\), and let \(\psi\) be a host character of parity
\(\psi(-1)=(-1)^a\).
Fix
\[
       \ell\notin\{2,3,7\},\qquad
                         \ell-1\nmid2a .
 \tag{SP1}
\]
Every odd-primary conductor component of a host character
is quadratic: the odd ramification comes from the two
quadratic radicals.
Consequently
\[
        \psi=\psi^{(\ell)}\widetilde\omega_\ell^{\,b},
       \qquad b\in\{0,(\ell-1)/2\},
 \tag{SP2}
\]
where \(\widetilde\omega_\ell\) is the global finite
Teichm\"uller character obtained from the cyclotomic character,
and \(\psi^{(\ell)}\) has conductor prime to \(\ell\).
For \(\ell\nmid N\), take \(b=0\).
For \(\ell=q\) or \(d\), the nonzero choice of \(b\) is
exactly the possible quadratic component of conductor \(\ell\).

Choose \(j\) with
\[
                       j\equiv a+b\pmod{\ell-1},
 \qquad
              \chi'=\psi^{(\ell)},\quad
              K=\mathbb Q(\sqrt{-3}),\quad\eta=\chi_K .
 \tag{SP3}
\]
Then
\[
 \chi'(-1)=(-1)^j,\qquad
  2j\equiv2a\pmod{\ell-1},\qquad
       \bar\varepsilon^{\,j}\bar\chi'
                       =\bar\varepsilon^{\,a}\bar\psi .
 \tag{SP4}
\]
The finite character
\(\varphi=\chi'\eta\) has conductor \(M>1\) prime to
\(\ell\).
It is nontrivial because \(\eta\) is ramified at \(3\),
whereas the sufficiently large principal host has no
ramification at \(3\).
This verifies the nontriviality assumption of the classical
seed construction, including the cases \(\ell=q,d\).

Let \(N_\ell\) be a prime-to-\(\ell\) common conductor of
\(\chi'\) and \(\eta\), for example
\(3N/\ell^{v_\ell(N)}\).
Choose a fresh prime
\[
                       w\equiv1\pmod{\ell N_\ell},
 \tag{SP5}
\]
and a nontrivial order-\(\ell\) Dirichlet character
\(\theta\) modulo \(w\).
Such a character exists because \(\ell\mid w-1\).
It has conductor \(w\), is even, and cuts out a totally
real cyclic field \(F_\theta/\mathbb Q\) of degree \(\ell\).
Equation (SP5) gives
\(\chi'(w)=\eta(w)=\varphi(w)=1\).
Thus all conductor, split-at-\(w\), parity, and nontriviality
requirements of the classical Eisenstein construction hold.
Its weights may be taken
\[
                           m\equiv j\pmod{\ell-1},
                         \qquad k=m+1
 \tag{SP6}
\]
arbitrarily large.

\subsection{The full local packet after the adjustment}
The representation constructed at weight \(m\) has global
residual character
\(\bar\varepsilon^{\,m}\bar\chi'\bar\eta\).
On the CM field containing \(K\) and \(F_\theta\), its
rank-three adjoint therefore has residual tuple
\[
        \bar\varepsilon^{\,a}\bar\psi,\quad1,\quad
                      (\bar\varepsilon^{\,a}\bar\psi)^{-1},
 \tag{SP7}
\]
by (SP4).
This is an equality of global residual characters.

The finite inertia types away from \(\ell\) agree
\emph{before reduction}.
Indeed \(\widetilde\omega_\ell\) is unramified at every
finite prime different from \(\ell\).
Consequently \(\psi|_{I_p}=\chi'|_{I_p}\) there, while
\(\eta,\theta\) vanish after the indicated field restriction.
The adjoint's type
\(\chi'\oplus1\oplus\chi'^{-1}\) is exactly the original
host finite type at those places.

At a place above \(\ell\), local class field theory identifies
the torsion of the local abelianization with the torsion
in the local multiplicative group.
The cyclotomic character on it takes values in
\(\mathbb Z_\ell^\times\).
It kills its \(\ell\)-primary torsion, because
\(\mathbb Z_\ell^\times\) has no nontrivial finite
\(\ell\)-power subgroup for odd \(\ell\).
The finite host characters have \(2\)-power order and
therefore kill the same \(\ell\)-primary torsion.
This is precisely the subgroup denoted
\srcid{GFrame.mu} in the source.

On the remaining finite torsion, the cyclotomic values are
Teichm\"uller units, so
\[
           \varepsilon^m
              =\varepsilon^a\widetilde\omega_\ell^{\,b}
 \tag{SP8}
\]
there, by (SP6).
The ordinary unit-root quotient of the seed is unramified
and kills these unit torsion elements.
Thus the ordered ordinary characters of its adjoint have
exactly the same finite-torsion restrictions as the original
packet.
Their principal-unit powers vary with the arithmetic
weight \(m\), as allowed in the common ordinary weight
algebra.
No equality between the distinct principal-unit powers
\(\varepsilon^m\) and \(\varepsilon^a\) is being claimed.
Equations (SP7)--(SP8), the fixed away-\(\ell\) types, and
the common polarization put the seed in the same prescribed
ordinary torsion component.
Applying the same fixed algebraic scalar normalization to
both packets preserves all these equalities.

\subsection{A finite-unramified CM tower with a scalar prime}
Choose a fresh prime \(p_0\), different from \(w\), with
\[
                 p_0\equiv-1\pmod{\ell N},\qquad
                              p_0\equiv1\pmod3 .
 \tag{SP9}
\]
The congruences are compatible because \(3\nmid\ell N\).
Dirichlet's theorem provides arbitrarily large such primes.
It splits in \(K\).
At this prime,
\[
 \bar\varepsilon(\operatorname{Frob}_{p_0})=-1,\qquad
 \bigl(\bar\varepsilon^{\,a}\bar\psi\bigr)
                      (\operatorname{Frob}_{p_0})=1.
 \tag{SP10}
\]
The second equation is the host parity:
\((-1)^a\psi(-1)=1\).
Thus all three residual adjoint eigencharacters coincide
at this initial Frobenius.

Let \(\mathcal P\) be any prescribed finite set of rational
primes at which the eventual CM extension must split
over its real part; include \(2,3,\ell\), and all primes
of the intended tame level.
The following construction also permits arbitrary prescribed
lower bounds for its real degree.

Choose a fresh positive prime \(D\) satisfying
\[
 \begin{gathered}
 D\equiv1\pmod4,\qquad D\equiv2\pmod3,\qquad
                      \left(\frac{3D}{p_0}\right)=1,\\
 \left(\frac Dp\right)=\left(\frac{-1}p\right)
 \quad\left(p\in\mathcal P,\ p\ge5,\
                     \left(\frac{-3}p\right)=-1\right).
 \end{gathered}
 \tag{SP11}
\]
It avoids \(\mathcal P\cup\{p_0,w\}\).
Each displayed Legendre condition is a nonempty choice of
unit residue classes at its respective odd prime.
There is no extra condition at \(p_0\) in the second line:
it splits in \(K\).
The Chinese remainder theorem and then Dirichlet's theorem
give a positive prime \(D\) in one of these simultaneous
unit classes.

Put
\[
            F_0^+=\mathbb Q(\sqrt{3D}),\qquad
                     L_0=F_0^+K.
\]
The third quadratic subfield of \(L_0\) is
\(\mathbb Q(\sqrt{-D})\).
Since \(D\equiv1\pmod4\), their discriminants are
\[
       -3,\qquad 12D,\qquad-4D .
\]
The discriminant formula for a biquadratic field gives
\[
        \operatorname{disc}(L_0)=144D^2
                          =\operatorname{disc}(F_0^+)^2 .
 \tag{SP12}
\]
The relative discriminant ideal therefore has norm one.
Thus \(L_0/F_0^+\) is unramified at \emph{every finite place}.

For \(p\in\mathcal P,\ p\ge5\), if \(p\) splits in \(K\),
then every place of \(F_0^+\) above it splits in \(L_0\).
If \(p\) is inert in \(K\), (SP11) gives
\[
                 \left(\frac{3D}p\right)
                    =\left(\frac{-3}p\right)=-1 .
\]
The two local quadratic fields \(K_p\) and \(F_{0,p}^+\)
are then the same unramified quadratic extension of
\(\mathbb Q_p\).
Their compositum splits over \(F_{0,p}^+\), so again the
CM extension splits at the desired place.
At \(3\), the ratio of the two quadratic radicands is
\((3D)/(-3)=-D\).
It is a square in \(\mathbb Q_3^\times\), since
\(D\equiv2\pmod3\).
Thus \(F_{0,3}^+=K_3\) and the CM extension splits there.

To arrange splitting at \(2\), choose another fresh positive
prime \(B\) with
\[
                      B\equiv5\pmod8,\qquad
                           \left(\frac B{p_0}\right)=1 .
 \tag{SP13}
\]
Again CRT and Dirichlet provide it, outside every already
specified finite set.
The field \(\mathbb Q_2(\sqrt B)\) is the unramified
quadratic extension of \(\mathbb Q_2\).
The real field \(F_0^+\) is ramified quadratic at \(2\),
since \(3D\equiv3\pmod4\).
The relative extension in (SP12) is unramified there,
and is killed by adjoining this unramified quadratic
field.  If it was already split there, splitting is
preserved.
Hence, after adjoining the real field \(\mathbb Q(\sqrt B)\),
the CM extension is split at \(2\) as well.
All previously obtained splitting conditions are preserved
by base change.

For each depth \(r\ge0\), let \(C_r\) be the unique
cyclic degree-\(\ell^r\) subfield of
\(\mathbb Q(\zeta_{\ell^{r+1}})\).
It is the field fixed by the Teichm\"uller subgroup of
order \(\ell-1\).
In particular \(C_r\) is totally real and totally ramified
at \(\ell\).
Its conductor is an \(\ell\)-power, whereas the conductors
of the preceding real fields are prime to \(\ell\).
This is the cyclotomic depth choice required for the
uniform class-group argument below.
Take
\[
 \begin{split}
 F^+&=F_0^+\mathbb Q(\sqrt B)F_\theta C_r,\\
                         L&=F^+K .
 \end{split}
 \tag{SP14}
\]
The fixed real biquadratic field has degree four.
The order-\(\ell\) field \(F_\theta\), of conductor \(w\),
is disjoint from the \(\ell\)-power conductor field \(C_r\):
their intersection would be an extension of \(\mathbb Q\)
unramified at every finite prime.
It is also disjoint from the biquadratic field.
Consequently
\[
                 [F^+:\mathbb Q]=4\ell^{r+1},
                        \qquad [L:F^+]=2 .
 \tag{SP15}
\]
The real degree is even.
At \(\ell\) the fixed real fields are unramified and
\(C_r\) is totally ramified of degree \(\ell^r\).
Every ordinary local degree is therefore at least
\(\ell^r\).
The global and all ordinary local degrees can exceed
any prescribed bound.
The uniform arithmetic bound for the three Selmer
dimensions denoted \(C_0\) is proved separately, before
this depth is chosen.
The fields are abelian over \(\mathbb Q\), \(L\) is CM,
and \(F^+\) is its maximal real subfield.
Also \(\eta|_{G_L}=\theta|_{G_L}=1\).

The finite-unramified property (SP12) survives every real
base change: an unramified local extension remains
unramified or split after finite base change.
It follows that \(L/F^+\) is unramified at every finite
place and split at every prime in \(\mathcal P\).
Thus it satisfies the exact \srcid{UnrSplit} condition.
The fixed fields in (SP14) have conductor prime to
\(\ell\); adjoining \(C_r\) gives ramification index
\(\ell^r\) at \(\ell\).
It follows that \(L\) is locally disjoint from the
tamely ramified extension
\(\mathbb Q_\ell(\mu_\ell)/\mathbb Q_\ell\),
whose degree is \(\ell-1\).
The image of \(\bar\varepsilon\) on inertia over \(L\)
is therefore still all of \(\mathbb F_\ell^\times\).
In particular no completion of \(L\) at \(\ell\)
contains \(\mu_\ell\) or nontrivial \(\ell\)-power roots
of unity.
Absolute unramifiedness at \(\ell\) is not needed.
Equations (SP1) and (SP4) show that
\[
               (\bar\varepsilon^{\,j}\bar\chi')^2
                                      \ne1\quad\hbox{on }G_L .
 \tag{SP16}
\]
This verifies the actual square-separation field condition,
not just the corresponding condition over \(\mathbb Q\).

Finally \(p_0\) splits in both real quadratic factors in
(SP14), by (SP11) and (SP13), and splits in \(K\).
It is unramified in \(F_\theta C_r\).
All the remaining residue degrees divide
\(\ell^{r+1}\), hence are odd.
For every \(u\mid p_0\) in \(L\), with residue degree \(f_u\),
restriction of unramified Frobenius gives
\[
 \bar\varepsilon(\operatorname{Frob}_u)
                       =(-1)^{f_u}=-1,\qquad
 (\bar\varepsilon^{\,a}\bar\psi)(\operatorname{Frob}_u)=1.
 \tag{SP17}
\]
Thus the scalar residual Frobenius and nontrivial
cyclotomic value survive the \emph{actual selected tower}.
The same conclusion holds for the adjusted character in
(SP4).
Any residual lattice need only have this scalar
characteristic polynomial; it need not have a scalar
matrix.  The matrix-minus-scalar is nilpotent after
reduction, which is exactly the local input used in the
separate \(N4\) matrix argument.
It also supplies the element with nontrivial
\(\bar\varepsilon\) required by first-Selmer trace detection.

\subsection{The definite rank-three group}
There is now no separate existence or parity problem for
the required definite unitary pure inner form.
Use the hermitian form
\[
                             h=1_3
\]
on \(L^3\).
It is positive definite at every real place of \(F^+\).
At a finite split place its unitary group is
\(\mathrm{GL}_3(F_v^+)\).
At a finite nonsplit place \(L_v/F_v^+\) is unramified.
Norms of units are surjective there, so \(-1\) is a norm.
The hermitian forms \(1_3\) and
\(\operatorname{diag}(1,-1,1)\) have the same determinant
class and hence are isometric over that local field.
The latter contains a hyperbolic plane.
Thus \(U(h)\) is quasi-split at every finite place.
Both forms are global hermitian forms over the same
quadratic extension, so their comparison defines a
pure inner twist, rather than merely an abstract inner
form.
The even degree in (SP15) is consistent with the global
product condition; it has been produced by the field
construction itself.

The constructed adjoint seed has a simple cuspidal
conjugate self-dual rank-three parameter.
With this group, the prescribed split finite local
members, the spherical members at the nonsplit finite
places, and the compact cohomological members at infinity
are the relevant local choices.
The separate unitary packet/vector realization must place
those choices in the ordinary module at its fixed type
and weight specialization.
The field construction and torsion matching above supply
its actual arithmetic hypotheses.

\paragraph{Correspondence and exact scope.}
The prime-to-\(\ell\) finite-character requirement is
\srcid{FieldPairCharacterTripleConditions.aux.cop} together
with \srcid{chi\_split} in
\srcid{IdeleGenerators.SlopeTestRemainingChoices}.
The local torsion subgroup is
\srcid{WeilGroupCharacters.UniversalOrdinaryDeterminantRing.GFrame.mu}.
The finite-unramified and split-at-level requirements are
\srcid{ModularityArgument.SymmetricPowerTransfer.UnrSplit}.
All these paths are under
\srcid{Development.AutomorphicArithmetic.AutomorphicVarieties}
at \srcid{alpoge-source-7cbf7a0c}.

This proves the packet adjustment at every coefficient
prime in (SP1), including \(\ell\mid N\), and gives explicit
common fields satisfying the seed's character, square,
finite-unramified, splitting, degree, and distinguished
Frobenius requirements.
It proves these assertions for the tower (SP14).
An independently chosen CM tower is not entitled to
(SP17) without an analogous residue-degree check.


\subsection{Inserting a stable cuspidal seed in the ordinary module}
\label{subsec:sr-stable-seed-insertion}
The local lattice and ordinary normalizations in
Section~\ref{subsec:sr-case-two-hecke} apply in arbitrary fixed rank.
Here is the precise specialization needed for the rank-three seed.
It uses the same finite function module at every auxiliary level.

\begin{srlemma}\label{lem:sr-stable-three-seed-insertion}
Let \(L/F\) be the actual CM tower, with \(F\) totally real,
\(L/F\) unramified at finite places and split at all places in
the prescribed finite set \(S\), including those above \(\ell\)
and the distinguished auxiliary place \(v_0\).
Suppose the definite hermitian group \(U(3)\) is a pure inner
form of its quasi-split unitary group and is quasi-split at
every finite place.
Let \(\Pi\) be a regular algebraic conjugate-self-dual
cuspidal representation of \(\mathrm{GL}_3(\mathbb A_L)\)
with the following actual local data:
\[
 \begin{gathered}
 r_{\ell,\iota}(\Pi)=r_0,\qquad
 r_0\bmod\mathfrak m\simeq
                  \bar\chi_0\oplus\bar\chi_1\oplus\bar\chi_2,\\
 \mathrm{WD}(r_0|_{G_{L_w}})^{F\text{-ss}}
       =\bigoplus_{i=0}^2 \lambda_{wi},\quad N=0
                    \quad(w\nmid\ell),\\
 \lambda_{wi}|_I=[\bar\chi_i]|_I
        \quad\hbox{at the prescribed split tame places},\\
 r_0|_{G_{L_w}}\ \hbox{ordinary with the prescribed ordered
 rank-one characters }\psi_{wi}
                    \quad(w\mid\ell).
 \end{gathered}
 \tag{S1}
\]
Assume at \(\ell\) that the local factor of \(\Pi\) is a
principal series with the ordinary refinement specified by
these characters. At nonsplit finite places outside \(S\),
assume it is unramified.
Assume also that \(\Pi\) is unramified at \(v_0\).
Then its seed point is supported by the residual ordinary
function module with these tame types and these weight
and ordinary eigenvalues.
In particular the kernel \(J_0\) of the actual ordinary
Hecke quotient vanishes at that point.
\end{srlemma}
\begin{proof}
Use
\href{https://arxiv.org/pdf/1409.3731v3}
{Kaletha--M\'inguez--Shin--White, Theorem 1.7.1}
only in the proved generic-parameter, pure-inner-twist
scope specified immediately after its statement.
The simple cuspidal parameter of \(\Pi\) has trivial global
component group modulo the center.
Its compact cohomological archimedean members therefore
occur with the prescribed split local members and
spherical members at every nonsplit finite place.
The hypotheses on the actual hermitian group supply
the pure-inner and global existence conditions.
This gives an actual automorphic representation \(\pi\)
of the chosen definite unitary group, with base change \(\Pi\).

At a split tame place use the coefficient representation
\[
                   \sigma^{\mathrm{crys}}(\tau_w)^\vee,
       \qquad \tau_w=\bigoplus_{i=0}^2[\bar\chi_i]|_I .
 \tag{S2}
\]
The local type criterion
\href{https://math.uchicago.edu/~emerton/pdffiles/bonn-lectures-survey-article.pdf}
{Emerton--Gee, Theorem 9.2.4}
and (S1) give
\(\operatorname{Hom}(\sigma^{\mathrm{crys}}(\tau_w),\Pi_w)
\ne0\).
The dual in (S2) is exactly the coefficient convention for
the function module. Choose its stable finite-free integral
lattice as in (H3).

At \(w\mid\ell\), the chosen principal-series refinement
has a one-dimensional Jacquet eigenspace.
After taking a sufficiently deep ordinary level, a
nonzero vector on this eigenspace is represented by a
local vector of the principal series.
This follows directly from the induced-function model:
on the open Bruhat cell prescribe its value at the identity,
extend by the inducing character and a sufficiently small
right compact open subgroup, and sum over the finite
unipotent quotient for the ordinary projector.
The normalized Jacquet action gives
\[
 \begin{aligned}
 \psi_{wi}(g)\epsilon_\ell(g)^i\,
                      \theta(e_{wi}(y))&=1
                  &&(\operatorname{Art}_w(g)=y\in\mathcal O_{L_w}^\times),\\
 U_{wj}&=\prod_{i=0}^j
                \psi_{wi}(g_w)\epsilon_\ell(g_w)^i
                  &&(\operatorname{Art}_w(g_w)=\varpi_w).
 \end{aligned}
 \tag{S3}
\]
These are (H5)--(H6) in rank three, so both the function
coefficient character \(\theta^{-1}\) and the operator
normalization agree with the common ordinary module.
The right-hand sides in the second line are units:
the \(\psi_{wi}\) are actual continuous Galois characters,
and the chosen local reciprocity lift has cyclotomic
value a unit.
The normalized ordinary projector thus acts as the
identity on the refined line.

Choose spherical vectors at the other unramified places.
At the distinguished auxiliary place use a spherical vector
inside the prescribed smaller level.
Tensor these with the local vectors above and the
cohomological archimedean factor.
The resulting automorphic vector is a nonzero classical
ordinary vector in precisely the finite function space used
to construct the common module.
All its coefficients, finite type matrices and eigenvalues
are defined over a finite coefficient extension.
There are finitely many double cosets.
Clear their finitely many denominators and scale to obtain
a primitive integral vector.
This operation retains every Hecke eigenvalue and
ordinary eigenvalue in (S3).

The spherical residual polynomial and (S3) are the
specified residual Hecke system.
The residual direct-summand projector therefore acts as
the identity on this vector. It survives in that
localized ordinary module.
The rank-independent coefficient-change and classical-control
isomorphisms identify its classical specialization with the
corresponding specialization of the same module over
\(\Lambda\).
Thus its Hecke eigensystem factors through the actual
closed Hecke algebra in \(\operatorname{End}_\Lambda(M)\).
The universal ordinary determinant maps to this algebra
with kernel \(J_0\). Evaluation at the constructed vector
annihilates \(J_0\), as claimed.
\end{proof}

For the rank-three adjoint seed, the preceding hypotheses are
checked on the actual seed, rather than imposed on an
arbitrary ordinary deformation.
Its construction over \(\mathbb Q\), cyclic CM base change
and Gelbart--Jacquet lift give the stable cuspidal \(\Pi\).
The local Galois representation after the inverse-cyclotomic
normalization is the actual \(r_0\) in (S1).
The local \(\ell\)-factor is unramified; its ordered Satake
parameters are those of the three crystalline ordinary
characters.
Their Jacquet refinement has the unit eigenvalues (S3).
The strict local adjoint root-ratio bound excludes linked
inducing characters, so this is an irreducible principal
series. The stated seed inertia at tame places gives (S2)
with \(N=0\).
The field construction gives a definite pure inner
\(U(3)\) quasi-split at finite places.
Consequently Lemma~\ref{lem:sr-stable-three-seed-insertion}
inserts this seed in the common module; no additional
existence statement for a modular ideal is used.

% Rank-independent part of THidaU, TTypeEx, TThetaEx,
% UnitaryHeckeData and ExistsHeckeIdeal, specialized to the
% stable rank-three cuspidal seed. The characteristic-zero
% eigenvector and residual support are in the SAME module.


\subsection{The actual rank-two companion seed on the same packet}
\label{subsec:generic-rank-two-companion}

The rank-three propagation uses the rank-two space
for the paired outer block.
Its initial point can be constructed on the same
fields and with the same auxiliary order-\(\ell\)
character as the rank-three seed.
We give that construction explicitly.

Retain
\(\Psi=\bar\varepsilon^j\bar\chi'
 =\bar\varepsilon^a\bar\psi\),
\(\eta=\chi_{\mathbb Q(\sqrt{-3})}\),
and \(\ell\notin\{2,3,7\}\).
Run the proved Eisenstein construction with
\[
                    \varphi_2=(\chi')^2\eta,
                         \qquad k_2=2m+1,\quad
                            m\equiv j\pmod{\ell-1}.           \tag{RS1}
\]
The finite character \(\varphi_2\) is primitive at
its own conductor, nontrivial, and of conductor
prime to \(\ell\).
Nontriviality is supplied by its factor at three.
It is odd, as required for the odd weight \(k_2\).
The same fresh \(w\) satisfies
\(\varphi_2(w)=1\), and the same \(\theta\)
has order \(\ell\), conductor \(w\), and trivial
restriction to \(G_L\).
Thus there are ordinary newforms \(f_2\) at all
sufficiently large weights (RS1).
Their actual Galois representations satisfy
\[
 \det\rho_{f_2}=\varepsilon^{2m}(\chi')^2\eta\theta,
 \qquad
 \bar\rho_{f_2}^{\rm ss}
                =1\oplus\bar\varepsilon^{2m}(\bar\chi')^2\eta .
                                                               \tag{RS2}
\]
At \(\ell\) they are crystalline ordinary of
weights \(2m,0\), with rank-one unramified
quotient; away from \(\ell\) their inertia type
is \(1\oplus(\chi')^2\eta\theta\), with \(N=0\).

These representations are irreducible on every
open subgroup, even in a case where
\(\Psi^4=1\).
Indeed failure of strong irreducibility in
dimension two gives either finite projective
image or a nontrivial quadratic self-twist
\(\alpha\).
Finite projective image contradicts the two
distinct Hodge--Tate weights on an open local
subgroup.
In the self-twist case, reduction of (RS2)
gives
\[
 \{1,\bar\Phi_2\}
          =\{\bar\alpha,\bar\alpha\bar\Phi_2\},
       \qquad\bar\Phi_2=\Psi^2\eta .
\]
Since \(\alpha\) is nontrivial quadratic in
odd characteristic, this forces
\(\bar\alpha=\bar\Phi_2\).

An isomorphism of the local ordinary
representation with its \(\alpha\)-twist
preserves its unique higher-weight stable line.
Uniqueness follows because any map from a
rank-one character of Hodge--Tate weight \(2m\)
to the rank-one weight-zero quotient is zero.
It consequently identifies the weight-zero
quotient with its \(\alpha\)-twist, which
forces \(\alpha|_{G_{\mathbb Q_\ell}}=1\).
But the restriction of
\(\bar\Phi_2=\Psi^2\eta\) to a place above
\(\ell\) of the selected real base is nontrivial:
\(\eta\) is trivial on that local group and
\(\Psi^2\ne1\) there.
This is a contradiction.
This argument uses the actual full local
non-equality and does not assume that
\(\Psi^4\ne1\) globally.

On \(G_L\), put
\[
             s_2=\rho_{f_2}\varepsilon^{-m}(\chi')^{-1}.
                                                               \tag{RS3}
\]
Its determinant is one.
Its residual ordered tuple is
\((\Psi,\Psi^{-1})\), its ordinary weights
are \(m,-m\), and its away-\(\ell\) inertia
type is \((\chi',(\chi')^{-1})\).
The \(\eta,\theta\) factors disappear on \(G_L\).
These are the same outer characters, in the
same order, as for the rank-three adjoint seed.

The pairing is symmetric conjugate self-dual,
rather than only an alternating self-dual
pairing.
Let \(A_2\) be the standard alternating
matrix on the two-dimensional space and use
the rational representation in (RS3)
before restriction.
For a complex conjugation \(c\),
\(\det s_2(c)=-1\) and \(s_2(c)^2=1\).
Hence
\[
                      P_2=A_2s_2(c),\qquad P_2^t=P_2 ,
\]
and, on \(G_L\),
\[
                    s_2(g)^tP_2s_2(cgc^{-1})=P_2 .
                                                               \tag{RS4}
\]
Both equations follow from
\(s_2(c)^tA_2s_2(c)=-A_2\) and
\(s_2(g)^tA_2s_2(g)=A_2\).
This supplies the exact paired rank-two
polarization.

For standard unitary normalization use a
fixed concrete CM character.
The elliptic curve
\[
                               E:\ y^2=x^3+1
\]
has CM by \(\mathbb Z[\zeta_3]\), with all
endomorphisms defined over
\(K=\mathbb Q(\sqrt{-3})\).
Over a finite coefficient extension its
Tate module restricted to \(G_K\) splits
into the characters \(\lambda,\lambda^c\).
The Weil pairing gives
\(\lambda\lambda^c=\varepsilon\).
Take
\[
                              \kappa=\lambda^{-1}|_{G_L}.
                                                               \tag{RS5}
\]
Then
\(\kappa\kappa^c=\varepsilon^{-1}\),
so \(s_2\kappa\) has the standard rank-two
unitary multiplier.
The displayed integral equation of \(E\)
has discriminant \(-432\), so all its
bad finite primes lie above \(2,3\),
already included in the level packet.
At the generic coefficient primes here it
has good reduction, and its characters are
crystalline.
Its finite inertia characters away from
\(\ell\) are killed by \(6\).
To see this directly, pass locally to a finite
extension over which \(E\) has good reduction.
The finite inertia action on the good special
fibre commutes with the specialized
\(\mathbb Z[\zeta_3]\)-action, since these
endomorphisms are defined over \(K\).
In the endomorphism algebra its centralizer is
\(K\); intersecting with integral automorphisms
gives \(\mathcal O_K^\times=\mu_6\).
This proves the assertion also at \(2,3\),
where the full good-fibre automorphism group
can be larger.
Since \(\ell\ge5\), this is prime-to-\(\ell\)
inertia.
These characters are retained in the common
residual and tame coefficient packet for the
scalar twist (RS5).
Thus (RS5) does not add an uncontrolled
finite ramification set.

The transfer identity in \srcid{IsTwist.ver}
also holds with its exact sign.
For \(t\in G_{\mathbb Q}\setminus G_K\), the
Tate representation interchanges the two CM
lines, so in their basis it has matrix
\(\left(\begin{smallmatrix}0&u\\v&0\end{smallmatrix}\right)\).
Its determinant is \(\varepsilon(t)=-uv\).
The square acts on either line by \(uv\);
hence
\[
        \lambda(t^2)=-\varepsilon(t),\qquad
        \kappa(t^2)=-\varepsilon(t)^{-1}.                     \tag{RS6}
\]
After restriction this is the same identity
for the selected CM packet and every allowed
element of its nontrivial real coset.


Strong irreducibility makes cyclic base
change of \(f_2\) to \(L\) cuspidal.
Classical CM theory identifies \(\lambda\)
with its algebraic Hecke character.
Twisting that base change by the algebraic
characters in (RS3)--(RS5) therefore gives
an actual regular algebraic conjugate
self-dual cuspidal rank-two representation
with Galois realization \(s_2\kappa\).

Its local inducing characters are unlinked.
For the unramified complementary components,
the same Gelbart--Jacquet bound
\(|t|<1/4\) used in the rank-three seed
bounds the ratio of the two roots by
\(q^{2|t|}<q\).
At a ramified primitive-level component the
ramified ratio and unitarity give \(t=0\).
Residue-degree base change preserves the
strict ratio inequality.
The scalar twists (RS3)--(RS5) cancel
from every ratio.

Finally \(h=I_2\) is the actual definite
hermitian form over \(L/F^+\).
It is positive at every real place.
At a nonsplit finite place the CM
extension is unramified, and \(-1\)
is a norm of a unit; thus \(I_2\)
is isometric there to the hyperbolic
hermitian plane.
At split places its group is
\(\mathrm{GL}_2\).
The even degree of \(F^+\) is supplied
by the field construction.
This is the required pure inner
unitary group.
The rank-independent stable-parameter
descent and local ordinary-vector
construction, with the exact
\(\theta,U_j\) normalizations already
proved for the common module, insert
(RS3)--(RS5) in its rank-two localized
ordinary factor.
The residual tuple and finite types
are exactly its prescribed tuple,
including the scalar \(\kappa\).
Thus the actual rank-two Hecke ideal
vanishes at this seed, at arbitrarily
large compatible \(m\).

This is the rank-two initial-point
producer called \srcid{Seed2Alt}.
Its classical constituent is the same
\srcid{BirthEG.propEGP\_of\_printsH}
construction, with finite character
\((\chi')^2\eta\).
The symmetric pairing and scalar
normalization above are the data used
by
\srcid{CurveDVRSpecialization.SeedPtOf}
and \srcid{PrLinkedEnd}.
The construction supplies a stable
rank-two seed; it does not replace
the paired rank-three support
argument by individually polarized
residual constituents.
